% book02.tex — Book II: Geometric algebra. Scaffolded.

\section{Book II --- Geometric Algebra}
\label{sec:book-II}

Book II contains the famous identities of geometric algebra, including
the distributivity of multiplication over addition (II.1) and the
law of cosines in its purely-geometric form (II.12, II.13). Definitions
introduce the gnomon construction used throughout.

\begin{claim}[Proposition II.1: Distributivity of multiplication]
\label{prop:II.1}
If there be two straight lines, and one of them be cut into any
number of segments whatever, the rectangle contained by the two
straight lines is equal to the rectangles contained by the
uncut straight line and each of the segments.
\end{claim}
\begin{evidence}[Proof sketch of II.1]
\label{ev:II.1}
Construction follows Heath. To be filled in fully in a later version.
\dependson{II.1}{I.31}
\dependson{II.1}{I.34}
\dependson{II.1}{cn:2}
\end{evidence}

\begin{claim}[Proposition II.12: Obtuse-triangle generalisation of Pythagoras]
\label{prop:II.12}
In obtuse-angled triangles the square on the side subtending the
obtuse angle is greater than the squares on the sides containing the
obtuse angle by twice the rectangle contained by one of the sides
about the obtuse angle (namely that on which the perpendicular falls)
and the straight line cut off outside by the perpendicular.
\end{claim}
\begin{evidence}[Proof sketch of II.12]
\label{ev:II.12}
Apply II.4 (square of a sum) and I.47.  To be filled in fully.
\dependson{II.12}{I.47}
\end{evidence}

% TODO: Propositions II.2 through II.14 (excluding II.12) — translations
% follow Heath; dependency edges to be added per the rrxiv convention.
